% Build targets:  ssrn.tex (working paper: title-page date, online appendix appended)
%                 el.tex   (Economics Letters submission: main text only)
%                 appendix.tex (stand-alone online appendix for the journal)
\documentclass[preprint,12pt,authoryear]{elsarticle}
\usepackage[T1]{fontenc}
\usepackage{lmodern,amsmath,booktabs,microtype,float}
\usepackage[hidelinks]{hyperref}
\biboptions{authoryear}
\journal{Economics Letters}
\newcommand{\authoremail}{tcanann@latech.edu}
% Replication package location. Uncomment and fill in when available:
% \newcommand{\DataURL}{https://github.com/<user>/<repo>}
% \newcommand{\DataDOI}{10.5281/zenodo.<number>}

\ifdefined\SSRN
\makeatletter
\def\ps@pprintTitle{\let\@oddhead\@empty\let\@evenhead\@empty
  \def\@oddfoot{\footnotesize\itshape Working paper. This version: \today\hfill Comments welcome.}\let\@evenfoot\@oddfoot}
\makeatother
\fi

\begin{document}
\begin{frontmatter}
\title{When disclosure dates and prices diverge: Evidence from the Funeral Rule}
\author{Taylor J. Canann}
\ead{\authoremail}
\affiliation{organization={Department of Economics and Finance, College of Business, Louisiana Tech University},
  addressline={504 W. Texas Ave.}, city={Ruston}, state={LA}, postcode={71270}, country={USA}}

\begin{abstract}
A new effective date on a mandated price list is often read as evidence of repricing. In 348 audited annual intervals of Greater Austin funeral-home price lists (2013--2024), headline prices changed in 87.2\% of intervals with a new date and 32.8\% with an unchanged date, and the gap survives removing changes that reverse. The date flags 83.1\% of price changes and 73.9\% of stable intervals, so date-based linear probability models recover 57\% of true effects when date errors are unrelated to the regressors. For price-list age, a regressor built from the dates, they are not.
\end{abstract}

\begin{keyword}
Price disclosure \sep Funeral Rule \sep Price rigidity \sep Misclassification \sep Measurement error
\JEL C81 \sep D18 \sep E31 \sep L11 \sep L84
\end{keyword}
\end{frontmatter}

\section{Introduction}
Many price-disclosure rules require the disclosed document to carry a date. Under the U.S. Federal Trade Commission's Funeral Rule, every funeral provider must offer a General Price List (GPL) stating the date from which its prices apply \citep{ftc_funeral_rule_guide}. Anyone who cannot observe prices continuously is tempted to read an unchanged date as unchanged prices. Consumer organizations that survey these lists carry prices forward when no new list arrives: the Kansas City survey of 2020 marks 25 of 113 entries as retained from 2018, and the Finger Lakes survey of 2025 reports a list dated 2019 \citep{fcagkc2020,fcafl2025}. Whether that inference is safe matters: the frequency and timing of price changes are central moments in the study of price rigidity \citep{bils_klenow_2004,nakamura_steinsson_2008}.

I test it with a validation design. For the same provider in consecutive annual surveys, I observe both the reported effective date and the recorded prices, so the date can be scored against the prices it describes. Posted-price documents have long been used to time price changes \citep{kashyap_1995,fisher_konieczny_2000}; the question here is whether a document's own date records that timing.

Three results follow. First, the date is informative but noisy: prices change in 87.2\% of intervals with a new date and 32.8\% with an unchanged one, a gap that survey noise does not explain. Second, the error concerns timing, not frequency: dates change in 64.9\% of intervals and prices in 68.1\%, yet the date misclassifies one interval in five. Third, for regressions that use a date change as the outcome, the error implies an attenuation factor of 0.570, but only when the date's errors are unrelated to the regressors. For price-list age, a regressor built from the dates, that condition fails.

\section{Data and measurement}
The Funeral Consumers Alliance of Central Texas (FCACTX) publishes an annual price survey of funeral establishments in the five-county Greater Austin area \citep{fcactx_cost_comparisons}. For each establishment it reports the GPL effective date and standardized prices. The 2013--2024 waves contain 622 establishment-year records for 71 establishments, linked by telephone number.

I track three headline prices separately: direct cremation, immediate burial including the provider's least expensive casket, and the non-declinable basic-services fee. An interval is an establishment observed in consecutive annual surveys with a reported date at both ends and at least two headline prices comparable at both ends. Repricing, $R$, equals one if a comparable price changes by more than \$5; a date change, $D$, equals one if the reported effective date differs. The sample has 348 intervals from 53 establishments.

All 45 intervals initially flagged as same-date repricing were re-read against the survey pages, removing five false flags, and a random recheck of final flags was added (online appendix). The date studied is the effective date as the survey records it, which is how such dates enter economic data. FCACTX constructs the two bundle prices from provider documents, so they can include components priced outside the GPL, such as the casket in the burial measure.

\section{Results}
\subsection{How well does the date measure repricing?}
Prices changed in 197 of 226 new-date intervals (87.2\%) and 40 of 122 unchanged-date intervals (32.8\%; Table~\ref{tab:signal}, first row). Treating repricing as the truth, the date has sensitivity 0.831 (it flags 197 of 237 repricing intervals) and specificity 0.739 (it leaves 82 of 111 stable intervals unflagged). In linear probability models of $R$ on $D$ with establishment-clustered standard errors, the coefficient is 0.529 (0.063) with year effects, 0.557 (0.075) with establishment and year effects, and 0.559 (0.076) adding prior list age; wild-cluster bootstrap $p$-values are below 0.001.

The missed changes are not rounding. The 40 same-date repricing intervals contain 49 changed prices with a median absolute change of \$150 (7.3\%); in 32 of the 40 the largest change is at least \$100. The gap survives requiring all three prices (87.4\% versus 34.0\%), dropping the burial measure (75.4\% versus 22.9\%), a 5\% threshold (69.9\% versus 19.7\%), and a worst-case assignment of missing price outcomes, which leaves a gap of 26.4 percentage points (online appendix).

Could the same-date changes be survey noise? They are more often decreases than new-date changes (18 of 49 versus 55 of 403) and more often reversed the next year (6 of 40 versus 12 of 315). Both patterns come mainly from the burial measure, whose casket price comes from a separately dated list: 12 of its 22 same-date changes are decreases and 5 of 19 reverse. Same-date changes in direct cremation look like repricing (20 of 26 are increases and 1 of 20 reverses), while the basic-services fee, a single GPL line, changes under an unchanged date in only 2 of 112 links. Dropping every change that reverses leaves 85.8\% versus 28.7\% ($J=0.578$). Survey records cannot show whether a list was edited without being re-dated or whether separately priced components changed; either way, the recorded date misstates when the prices consumers compare changed.

\subsection{What the error does to regressions}
Suppose a researcher observes $D$ rather than $R$ and regresses it on covariates $X$. If the date's error rates do not depend on $X$, then
\begin{equation}
E[D\mid X] = (1-\text{spec}) + J\,E[R\mid X], \qquad J = \text{sens}+\text{spec}-1,
\end{equation}
so every linear-probability coefficient is scaled by $J$ \citep{hausman_abrevaya_scottmorton_1998}. In Austin $J=0.570$ (block-bootstrap 95\% interval 0.454--0.681): a date-based study recovers 57\% of a true effect.

The condition fails by construction for covariates computed from the dates. Prior list age, the time since the previous effective date, is the natural state variable in a hazard model of repricing; Table~\ref{tab:age} regresses each outcome on it. With establishment and year effects, a year of list age raises the probability of a date change by 0.088 (0.038) and of a price change by 0.045 (0.022). Nondifferential error would imply a date coefficient of $J\times0.045=0.025$; the estimate is 3.5 times that, and the block-bootstrap interval for the difference (0.006--0.138) excludes zero. Without establishment effects, the date coefficient reverses sign ($-0.076$), the price coefficient is near zero ($-0.008$), and age predicts a date change even holding the price change fixed ($-0.072$, wild $p=0.016$).

The date coefficients are largely mechanical. List age resets whenever the date changes, so with about seven intervals per establishment the within estimate inherits a \citet{nickell_1981}-type bias, while the pooled estimate mixes providers that rarely re-date with providers that re-date often. Simulating a constant hazard, with no duration dependence, on Austin's panel dimensions yields within coefficients of 0.063--0.091 and pooled coefficients of $-0.101$ to 0.003, depending on provider heterogeneity (online appendix). A date-based hazard regression can therefore report duration dependence, of either sign, that the date itself produces \citep[cf.][]{alvarez_borovickova_shimer_2025}.

\subsection{A second market}
Kansas City offers an out-of-sample check. FCA of Greater Kansas City publishes month-level GPL dates and prices at two- and three-year intervals \citep{fcagkc2018,fcagkc2020,fcagkc2023}. For direct cremation, excluding records marked as carried forward, a new GPL month again accompanies repricing far more often (62.7\% and 79.4\%) than an unchanged month (4.5\% and 14.3\%, from only 22 and 7 intervals). The error runs the other way, toward new dates without a changed price, so $J$ is 0.462 and 0.282. Matching the measure and gap in Austin gives $J=0.303$ for direct cremation over two-year links, inside that range, against 0.570 for annual links on any of the three prices (Table~\ref{tab:signal}).
The attenuation factor therefore depends on the price measure and interval and must be estimated for the data at hand.

\section{Conclusion}
A disclosed effective date is a useful but imperfect record of repricing. In Greater Austin it gets the frequency of price changes nearly right and the timing wrong in one interval in five, including changes of \$100 or more. It is most reliable for a single GPL line: an unchanged date accompanied an unchanged basic-services fee in 110 of 112 links. It is least reliable for the arrangement prices consumers compare. Studies that must rely on dates should estimate sensitivity and specificity for their own measure and interval, and rescale by $J$ only for regressors plausibly unrelated to the date's errors; for date-derived covariates such as list age, no rescaling applies. For disclosure policy, the FTC's consideration of online price-list posting \citep{ftc_anpr_2022} raises the value of dates that track prices.

\begin{table}[t]
\centering\footnotesize
\caption{Effective-date changes and observed repricing.}\label{tab:signal}
\setlength{\tabcolsep}{2.5pt}\scriptsize
\begin{tabular*}{\textwidth}{@{\extracolsep{\fill}}lccccccc@{}}
\toprule
 & Gap & & \multicolumn{2}{c}{Repriced/intervals (\%)} & & & \\
\cmidrule(lr){4-5}
Market, price measure & (yrs) & $N$ & New date & Same date & Sens. & Spec. & $J$ \\
\midrule
Austin, any of three & 1 & 348 & 197/226 (87.2) & 40/122 (32.8) & 0.831 & 0.739 & 0.570 \\
Austin, any of three & 2 & 325 & 255/271 (94.1) & 20/54 (37.0) & 0.927 & 0.680 & 0.607 \\
Austin, cremation & 2 & 355 & 233/299 (77.9) & 17/56 (30.4) & 0.932 & 0.371 & 0.303 \\
Kansas City, cremation, 2018--20$^{a}$ & 2 & 81 & 37/59 (62.7) & 1/22 (4.5) & 0.974 & 0.488 & 0.462 \\
Kansas City, cremation, 2020--23$^{a}$ & 3 & 75 & 54/68 (79.4) & 1/7 (14.3) & 0.982 & 0.300 & 0.282 \\
\bottomrule
\end{tabular*}
\begin{minipage}{\textwidth}\vspace{3pt}\scriptsize
Notes: Repricing: a comparable price changes by more than \$5. Sens.: share of repricing intervals with a new date; Spec.: share of stable intervals with an unchanged date; $J$ = Sens. + Spec. $-$ 1. ``Any of three'': direct cremation, immediate burial incl.~casket, or basic-services fee; at least two comparable. Cremation rows use all links at which direct cremation is comparable. $^{a}$Month-level dates from published survey tables; carried-forward records excluded.
\end{minipage}
\end{table}

\begin{table}[t]
\centering\footnotesize
\caption{Prior price-list age and the two measures of repricing, Greater Austin.}\label{tab:age}
\begin{tabular*}{\textwidth}{@{\extracolsep{\fill}}lcccc@{}}
\toprule
Outcome: & Date change & Price change & Date change & Price change \\
Also controls for: & -- & -- & Price change & Date change \\
 & (1) & (2) & (3) & (4) \\
\midrule
\multicolumn{5}{@{}l}{\emph{Establishment and year effects}}\\
Prior list age (years) & 0.088 & 0.045 & 0.065 & $-$0.005 \\
 & (0.038) & (0.022) & (0.037) & (0.027) \\
 & [0.016] & [0.062] & [0.096] & [0.869] \\[2pt]
\multicolumn{5}{@{}l}{\emph{Year effects only}}\\
Prior list age (years) & $-$0.076 & $-$0.008 & $-$0.072 & 0.034 \\
 & (0.022) & (0.018) & (0.021) & (0.017) \\
 & [0.006] & [0.673] & [0.016] & [0.099] \\[2pt]
\bottomrule
\end{tabular*}
\begin{minipage}{\textwidth}\vspace{3pt}\scriptsize
Notes: Linear probability models; 348 annual intervals, 53 establishments. Prior list age is the age of the effective date on the preceding survey's list. Standard errors clustered by establishment in parentheses; wild-cluster bootstrap $p$-values (Rademacher, 4,999 draws) in brackets. Under nondifferential date error the coefficient in column (3) is zero and column (1) equals $J\times$ column (2). Column (4) conditions on an outcome.
\end{minipage}
\end{table}


\section*{Data availability}
The replication package contains the FCACTX survey publications, the audited Austin panel, the Kansas City survey-derived files and code that reproduces every table and number. FCACTX permits public redistribution; FCA-GKC has authorized publication and redistribution. \ifdefined\DataDOI The package is archived at \url{https://doi.org/\DataDOI}\ifdefined\DataURL\ and maintained at \url{\DataURL}\fi.\else\ifdefined\DataURL The package is available at \url{\DataURL}.\else The package will be deposited in a public repository with a DOI; until then it is available from the author.\fi\fi

\section*{Declaration of generative AI and AI-assisted technologies in the manuscript preparation process}
During the preparation of this work the author used Claude (Anthropic) in order to re-run and extend the statistical code and to draft and edit the text. After using this tool/service, the author reviewed and edited the content as needed and takes full responsibility for the content of the published article.

\section*{Acknowledgments}
I thank the Funeral Consumers Alliance of Central Texas, the Funeral Consumers Alliance of Greater Kansas City and the Funeral Consumers Alliance of the Finger Lakes for their price surveys.

\bibliographystyle{elsarticle-harv}
\bibliography{references}

\ifdefined\SSRN
\clearpage
% Online appendix body. Generated by write_appendix_body.py from appendix_template.tex and output/appendix.json; edit the template.
\section*{Online appendix}
\setcounter{table}{0}
\renewcommand{\thetable}{A\arabic{table}}
\setcounter{section}{0}
\renewcommand{\thesection}{A.\arabic{section}}

\section{Data, audit and sample construction}
The Greater Austin panel is read from the FCACTX survey publications for 2013--2024 (there was no 2025 survey; the 2026 wave is not used). Prices and effective dates are taken from the published tables with layout-preserving rules, and unresolved values stay missing. Establishments are linked across waves by telephone number. A two-year link compares survey $t-2$ with survey $t$ for the same establishment, whether or not $t-1$ is observed, using the same rules as the annual links.

The 45 intervals initially flagged as same-date repricing were checked against the current- and preceding-year survey pages. Four price-row alignment errors and one date-parsing error removed five false flags; one real change became ineligible under the two-comparable-measure rule; and correcting a preceding date moved one adjacent interval into the same-date category, giving $45-4-1-1+1=40$. Ten randomly selected final flags (seed 20260907) and two targeted cases were then rechecked. All twelve kept their flag; one required a price-level correction (Gabriels Funeral Chapel, 2014). The audit verifies consistency with the survey pages, not with every underlying provider document.

Of 435 date-defined consecutive intervals, 87 lack two comparable headline prices. Setting every missing new-date outcome to no repricing and every missing same-date outcome to repricing gives a lower bound on the gap of $197/300-53/135=0.264$. The bound does not address records without a usable date.

Kansas City comparisons use FCA-GKC's 2018, 2020 and 2023 publications, which report month-level GPL dates. Branches are linked across publications, and records that the 2020 publication marks as retained from 2018 are excluded at either endpoint. Because dates are month--year fields, an unchanged month is not necessarily an unchanged day, and inclusion of the crematory fee in the 2018 direct-cremation measure is not explicitly established.

\section{The date signal by price measure and observation gap}
Table~\ref{tab:A1} scores the date against each price separately. For the basic-services fee, a single GPL line, the date is almost perfectly sensitive (0.980): the fee changed without a new date in 2 of 112 links. Its specificity is low (0.493) because most new lists leave the fee unchanged. The two bundle prices change under an unchanged date more often (28 of 127 for direct cremation and 22 of 119 for burial). Over two-year links, few dates stay unchanged, so sensitivity rises and specificity falls. The date has also become more informative over time: $J$ is 0.492 for intervals ending in 2014--2019 and 0.680 for 2020--2024.

\begin{table}[H]\centering\footnotesize
\caption{Date signal by price measure and observation gap.}\label{tab:A1}
\setlength{\tabcolsep}{2.5pt}\scriptsize
\begin{tabular*}{\textwidth}{@{\extracolsep{\fill}}p{4.6cm}cccccc@{}}\toprule
Price measure, links & $N$ & New date & Same date & Sens. & Spec. & $J$ \\\midrule
Basic-services fee, annual & 323 & 98/211 (46.4) & 2/112 (1.8) & 0.980 & 0.493 & 0.473 \\
Direct cremation, annual & 379 & 165/252 (65.5) & 28/127 (22.0) & 0.855 & 0.532 & 0.387 \\
Burial incl.~casket, annual & 341 & 159/222 (71.6) & 22/119 (18.5) & 0.878 & 0.606 & 0.485 \\
Any of three (main sample), annual & 348 & 197/226 (87.2) & 40/122 (32.8) & 0.831 & 0.739 & 0.570 \\
Direct cremation, two-year links & 355 & 233/299 (77.9) & 17/56 (30.4) & 0.932 & 0.371 & 0.303 \\
Any of three, two-year links & 325 & 255/271 (94.1) & 20/54 (37.0) & 0.927 & 0.680 & 0.607 \\
Kansas City, cremation, 2018--20 & 81 & 37/59 (62.7) & 1/22 (4.5) & 0.974 & 0.488 & 0.462 \\
Kansas City, cremation, 2020--23 & 75 & 54/68 (79.4) & 1/7 (14.3) & 0.982 & 0.300 & 0.282 \\
\bottomrule\end{tabular*}
\begin{minipage}{\textwidth}\vspace{3pt}\scriptsize Notes: Repriced/links (\%) for new-date and same-date links. Single-measure rows use every link at which that price is comparable, so samples differ from the main sample. Two-year links compare survey $t-2$ with $t$. Kansas City rows use month-level dates from published FCA-GKC survey tables and exclude records marked as carried forward.\end{minipage}
\end{table}


\section{Same-date changes and survey noise}
Table~\ref{tab:A2} compares same-date and new-date price changes on two symptoms of recording noise: decreases, and reversal to the earlier price in the following year. Same-date changes show more of both, and almost all of the excess is in the burial measure, whose least-expensive casket is priced on a separately dated casket price list. Removing every reversed change leaves the result in the main text. If changes whose next-year price is unobserved are also dropped, the rates are 149/181 (82.3\%) and 27/114 (23.7\%), with $J=0.578$. Same-date repricing occurs at 16 establishments, and none accounts for more than 15\% of the 40 intervals.

\begin{table}[H]\centering\footnotesize
\caption{Do same-date price changes look like survey noise?}\label{tab:A2}
\setlength{\tabcolsep}{3pt}
\begin{tabular*}{\textwidth}{@{\extracolsep{\fill}}p{3.6cm}cccccc@{}}\toprule
 & \multicolumn{3}{c}{Same date} & \multicolumn{3}{c}{New date}\\\cmidrule(lr){2-4}\cmidrule(lr){5-7}
Component & Changes & Decreases & Reversed & Changes & Decreases & Reversed \\\midrule
Direct cremation & 26 & 6 & 1/20 & 153 & 22 & 4/125 \\
Burial incl.~casket & 22 & 12 & 5/19 & 159 & 22 & 5/124 \\
Basic-services fee & 1 & 0 & 0/1 & 91 & 11 & 3/66 \\
All & 49 & 18 & 6/40 & 403 & 55 & 12/315 \\
\bottomrule\end{tabular*}
\begin{minipage}{\textwidth}\vspace{3pt}\scriptsize Notes: Changed component events (more than \$5) in the 348 annual intervals. Reversed: the price at $t+1$ is within \$5 of its value at $t-1$, among events whose next-year price is observed. Fisher exact $p$-values for same- versus new-date differences: decreases $<0.001$, reversals 0.009.\end{minipage}
\end{table}


\section{Robustness of the repricing gap}
Table~\ref{tab:A3} collects alternative samples, thresholds and exclusions. Except for the worst-case missing-data bound (0.264), the gap between new-date and same-date repricing stays between 0.502 and 0.652, and no single brand or establishment drives it.

\begin{table}[H]\centering\footnotesize
\caption{Robustness of the same-date versus new-date repricing gap, Greater Austin.}\label{tab:A3}
\begin{tabular*}{\textwidth}{@{\extracolsep{\fill}}p{6.2cm}ccc@{}}\toprule
Check & New date & Same date & Gap \\\midrule
Baseline & 197/226 (87.2) & 40/122 (32.8) & 0.544 \\
All three prices observed & 167/191 (87.4) & 36/106 (34.0) & 0.535 \\
Cremation or fee only (burial dropped) & 147/195 (75.4) & 25/109 (22.9) & 0.524 \\
Both surveys 2019 or later & 74/82 (90.2) & 13/52 (25.0) & 0.652 \\
Omit 2016--2017 survey redesign & 178/203 (87.7) & 34/110 (30.9) & 0.568 \\
Omit website/phone-sourced flags & 197/226 (87.2) & 37/119 (31.1) & 0.561 \\
Omit two audited histories & 189/217 (87.1) & 34/109 (31.2) & 0.559 \\
Threshold \$25 & 194/226 (85.8) & 40/122 (32.8) & 0.531 \\
Threshold 2.5\% & 182/226 (80.5) & 35/122 (28.7) & 0.518 \\
Threshold 5\% & 158/226 (69.9) & 24/122 (19.7) & 0.502 \\
Drop changes reversed next year & 194/226 (85.8) & 35/122 (28.7) & 0.572 \\
Worst-case missing price outcomes & & & 0.264 \\
Leave one coded brand out (range) & & & 0.524--0.560 \\
Leave one establishment out (range) & & & 0.523--0.572 \\
\bottomrule\end{tabular*}
\begin{minipage}{\textwidth}\vspace{3pt}\scriptsize Notes: Repriced/intervals (\%). Rows other than the last four and the reversal row reproduce the replication package's robustness outputs. Worst case: every missing new-date outcome set to no repricing and every missing same-date outcome to repricing. ``Drop changes reversed next year'' keeps changes whose next-year price is unobserved.\end{minipage}
\end{table}


\section{Prior list age: additional specifications and simulation}
Table~\ref{tab:A4} splits the price regression by date status and repeats both regressions on the sample at which direct cremation and the fee are comparable, where the date and price gradients are close. Table~\ref{tab:A5} reports the simulation described in the main text. With establishment effects, list age predicts a date change even though the simulated hazard is constant, because age resets when the date changes and the within transformation uses each provider's own average date-change rate. Without establishment effects, providers that rarely re-date both hold older lists and re-date less, which produces a negative coefficient.

\begin{table}[H]\centering\footnotesize
\caption{Prior list age: additional specifications.}\label{tab:A4}
\begin{tabular*}{\textwidth}{@{\extracolsep{\fill}}p{7cm}cc@{}}\toprule
Outcome and sample & Coefficient on prior list age & $N$ \\\midrule
Price change, same-date intervals only & 0.015 (0.043) [0.826] & 122 \\
Price change, new-date intervals only & $-$0.007 (0.025) [0.755] & 226 \\
Date change, cremation/fee sample & 0.082 (0.041) [0.049] & 304 \\
Cremation or fee change, cremation/fee sample & 0.065 (0.037) [0.076] & 304 \\
Date minus price change, establishment and year effects & 0.044 (0.039) [0.285] & 348 \\
Date minus price change, year effects only & $-$0.069 (0.023) [0.030] & 348 \\
\bottomrule\end{tabular*}
\begin{minipage}{\textwidth}\vspace{3pt}\scriptsize Notes: Linear probability models. The first five rows include establishment and year effects. Establishment-clustered standard errors in parentheses; wild-cluster bootstrap $p$-values (4,999 Rademacher draws) in brackets. The cremation/fee sample is the 304 intervals at which both direct cremation and the basic-services fee are comparable. ``Date minus price'' regresses $D-R$ on age and tests whether the two gradients are equal.\end{minipage}
\end{table}

\begin{table}[H]\centering\footnotesize
\caption{List-age coefficients when the true hazard is constant.}\label{tab:A5}
\begin{tabular*}{\textwidth}{@{\extracolsep{\fill}}lcc@{}}\toprule
Simulated process & Establishment and year effects & Year effects only \\\midrule
Homogeneous hazard, $p=0.65$ & 0.091 (0.027) & 0.003 (0.024) \\
Heterogeneous, Beta(2.6, 1.4) & 0.076 (0.023) & $-$0.068 (0.021) \\
More heterogeneous, Beta(1.3, 0.7) & 0.063 (0.020) & $-$0.101 (0.015) \\
Austin estimate, Table 2 column (1) & 0.088 & $-$0.076 \\
\bottomrule\end{tabular*}
\begin{minipage}{\textwidth}\vspace{3pt}\scriptsize Notes: Mean (standard deviation) of the coefficient on prior list age in a regression of the date-change indicator, over 500 simulated panels of 53 providers and seven annual intervals. Each provider re-dates with a constant probability drawn as shown (mean 0.65, matching Austin's 64.9\% date-change frequency), so there is no duration dependence. Initial list age is uniform on $[0,3]$ years; a list issued during an interval is uniform on $[0,1]$ years old at the next survey; otherwise age rises by one year.\end{minipage}
\end{table}


\fi
\end{document}
